%% Appendix D — RTDB Schema Reference
\chapter{Realtime Database Schema Reference}
\label{app:rtdb_schema}

\section{Full RTDB Tree}

\begin{lstlisting}[style=jsonlstyle,caption={Complete RTDB Structure}]
{
  "mms": {
    "nodes": {
      "1": {
        "state":            "idle",
        "roll_number":      null,
        "session_start":    null,
        "firmware_version": "2.0.0",
        "ip_address":       "192.168.0.101",
        "last_seen":        1718600000,
        "rssi":             -52
      },
      "2": {
        "state":            "offline",
        "firmware_version": "2.0.0",
        "ip_address":       null,
        "last_seen":        1718500000,
        "rssi":             null
      },
      "3": {
        "state":            "active",
        "roll_number":      "220002123",
        "session_start":    1718600000,
        "firmware_version": "2.0.0",
        "ip_address":       "192.168.0.103",
        "last_seen":        1718600060,
        "rssi":             -67
      }
    },
    "commands": {
      "3": {
        "command":      "stop",
        "issued_by":    "admin_firebase_uid",
        "timestamp":    1718600100,
        "acknowledged": false
      }
    },
    "ota": {
      "firmware_url":   "https://firebasestorage.googleapis.com/.../firmware/2.1.0.bin",
      "target_version": "2.1.0",
      "released_at":    1718600200
    },
    "bridge": {
      "status": {
        "state": "online",
        "ts":    1718600000
      }
    }
  }
}
\end{lstlisting}

\section{Node State Values}

\begin{table}[H]
\centering
\caption{Valid values for \texttt{/mms/nodes/\{id\}/state}}
\begin{tabular}{L{4cm}L{10cm}}
\toprule
\textbf{State} & \textbf{Description} \\
\midrule
\texttt{"offline"} & Node is not connected to MQTT. Published by broker's LWT mechanism. \\
\texttt{"idle"} & Node is online, no card in slot. \\
\texttt{"card\_present"} & Card detected in slot; RFID read in progress. \\
\texttt{"active"} & Session in progress; machine powered via relay. \\
\texttt{"error"} & Node encountered a hardware error and is in error recovery mode. \\
\texttt{"ota"} & Node is performing an OTA firmware update. \\
\bottomrule
\end{tabular}
\end{table}

\section{Command Lifecycle}

\begin{enumerate}[leftmargin=1.5em]
    \item Admin writes command to \texttt{/mms/commands/\{id\}}: \texttt{\{command: "stop", acknowledged: false\}}
    \item Bridge forwards command to MQTT
    \item Node receives, executes (or discards if stale)
    \item Node publishes session event
    \item Bridge or node sets \texttt{\{acknowledged: true, ack\_at: now\}}
    \item Bridge periodic sweep cleans up old acknowledged commands after 10 minutes
\end{enumerate}

\section{RTDB Security Rules}

\begin{lstlisting}[caption={v2/database.rules.json}]
{
  "rules": {
    "mms": {
      "nodes": {
        ".read": "auth != null",
        ".write": false
      },
      "commands": {
        "$machine_id": {
          ".read": "auth != null",
          ".write": "auth != null"
        }
      },
      "ota": {
        ".read": "auth != null",
        ".write": "auth != null"
      },
      "bridge": {
        ".read": "auth != null",
        ".write": false
      }
    }
  }
}
\end{lstlisting}

\begin{notebox}
RTDB rules apply only to client-side SDK connections. The bridge.py uses the Firebase Admin SDK, which bypasses all RTDB security rules. This is by design: the bridge must be able to write node states, which clients cannot write directly.
\end{notebox}
